% P10-HITL, 15.09.2026: Entscheidungszweck und belegte Zustandswirkung präzisiert.
% P10-4a, 14.09.2026: Mechanismus, Quellenreichweite und v4-Nachweis angeglichen.
% M08, 11.09.2026: Governance als Hauptstelle; Begründungspflicht
% pfadspezifisch (IngestionReviewAdapter vs. DecisionResolutionReviewAdapter).
\section{Governance und kontrollierte Mutation}
\label{sec:governance-mutation}

Ein vorbereiteter Änderungsvorschlag ist noch kein geltender
Projektinhalt. Für die Fortschreibung des vorhandenen Bestands
verbindet der kontrollierte Mutationspfad deshalb fachliche Prüfung,
menschliche Autorisierung, Anwendung und Speicherung. Er konkretisiert
Anforderung A6. Der initiale Seed schreibt dagegen ohne
Einzelautorisierung der entstehenden Elemente; dort ist das
Gestaltungsziel nur teilweise realisiert
(Abschnitt~\ref{sec:endtoend-pfade}).
Abbildung~\ref{fig:kontrollierter-mutationspfad} zeigt den Ablauf
am Beispiel der Anforderungseinordnung.

An den menschlichen Entscheidungspunkten (\emph{Human Gates}) entscheidet
der Mensch, welche fachliche Wirkung ein Vorschlag erhalten soll. Das ist auch bei
einem plausiblen Vorschlag erforderlich: Eine neue Aussage kann eine
bisherige Festlegung infrage stellen, ohne bereits die Entscheidung
über ihre Ablösung zu enthalten. HITL ermöglicht somit sowohl die
Korrektur von Vorschlägen als auch eine verbindliche fachliche Auswahl.

\begin{figure}[htbp]
  \centering
  \includegraphics[width=\textwidth]{figures/05/Abbildung5.6.pdf}
  \caption[Kontrollierter Mutationspfad]{Kontrollierter Mutationspfad
    am Beispiel der Anforderungseinordnung
    (H: menschlicher Entscheidungspunkt; eigene Darstellung).}
  \label{fig:kontrollierter-mutationspfad}
\end{figure}

\emph{Prüfung und Entscheidung.} Bei der Anforderungseinordnung
vergleicht der Agent den neuen Eingang mit dem vorhandenen Core und
erstellt einen typisierten Operationsplan. Eine deterministische
Prüfung kontrolliert unter anderem die zulässigen Operationstypen
und ihre Zielreferenzen. Bei reparierbaren Verstößen erhält der Agent
die Befunde zur Überarbeitung; die Versuchsschranke begrenzt diese
Rückkopplung (Abschnitt~\ref{sec:saeule-semantische-transformation}).
Nach bestandener Prüfung sieht der Mensch am Human Gate die
vorgeschlagene Operation, den eingehenden Inhalt, die Begründung und
gegebenenfalls den betroffenen Bestandstext. Er übernimmt einzelne
Operationen oder lehnt sie mit Begründung ab. Die Anwendung übernimmt
ausschließlich die freigegebene Menge. Sie konstruiert den nächsten Core-Zustand mit
den jeweiligen Elementen, Herkunftsangaben und Beziehungen
(Abschnitt~\ref{sec:core-projektfortschreibung}). Welche Felder
und Verbindungen entstehen, bestimmt die konkrete Operation;
die gemeinsame Speicheroperation erzeugt sie nicht selbst.

\emph{Speicherung.} Die resultierende Änderung gelangt zur
gemeinsamen Speicheroperation, der Speichernaht. Sie prüft
strukturelle Invarianten vor dem Schreiben und erhält bei einer
Änderung den bisherigen Bestand in der Historie. Die als Fehler eingestuften
Verstöße verhindern das Schreiben; weitere Auffälligkeiten werden
als Warnungen protokolliert. Geprüft wird beispielsweise, ob eine
\emph{deckt}-Beziehung von einem PBI auf eine vorhandene Anforderung
oder ein Architekturelement zeigt. Externe Herkunftsziele werden
dabei nicht vollständig auf ihre Auflösbarkeit geprüft. Die
Realisierung dieses Wächters, \emph{CoreKangal}, erläutert
Abschnitt~\ref{sec:realisierung-persistenz-aussenkopplung}. Die Speichernaht
erzwingt damit strukturelle Zulässigkeit, keine fachliche
Autorisierung: Diese muss im vorgelagerten Verarbeitungspfad
erfolgen. Auch wartende Vorschläge können bereits als begleitende
Governance-Information gespeichert werden, ohne dadurch zu
geltenden Anforderungen zu werden.

Drei Kontrollarten beantworten somit unterschiedliche Fragen
(Tabelle~\ref{tab:kontrollarten}). Insbesondere autorisiert die
Zustimmung zu einem Werkzeugaufruf noch nicht die fachlichen
Änderungen, die dieser Aufruf erst vorschlägt.

\begin{table}[htbp]
  \centering
  \small
  \begin{tabular}{p{0.26\textwidth}p{0.30\textwidth}p{0.32\textwidth}}
    \hline
    \textbf{Kontrollart} & \textbf{Leitfrage} & \textbf{Wirkort} \\
    \hline
    Werkzeugzustimmung (Tool Approval) & Darf diese Handlung
      ausgeführt werden? & agentische Bedienebene, für entsprechend
      als wirkungs- oder kostenträchtig ausgewiesene Handlungen \\
    fachliche Autorisierung (Human Gate) & Soll die fachliche
      Änderung gelten? & vor der Anwendung einer vorgeschlagenen
      Änderung, je Element \\
    Invariantenprüfung & Ist die resultierende Mutation strukturell
      zulässig? & gemeinsame Speichernaht, vor der Persistenz \\
    \hline
  \end{tabular}
  \caption{Drei Kontrollarten der Governance und ihre getrennten
    Leitfragen.}
  \label{tab:kontrollarten}
\end{table}

Im Besuchsfall F1 wird die fachliche Entscheidung an der Übernahme
neuer Informationen sichtbar: Der Einordnungsplan legte zwölf
Operationen vor, von denen der Mensch acht übernahm und vier mit
Begründung abwies. Zu den übernommenen Inhalten gehören die beiden
Anforderungen zur Besuchsankündigung und zur Übersicht für Pflegende.
Abgewiesen wurde beispielsweise eine vom Agenten als Verfeinerung
vorgeschlagene Änderung an REQ-42. Der Mensch bewertete sie als bloße
Bestätigung des Bestands; der vorgeschlagene Ersatztext hätte die
vorhandenen Angaben zu einnahmebezogenen Bemerkungen verdrängt.
Die protokollierte Auswahl bestimmt, welche Änderungen angewendet
werden. Die Prüfung von Vorlage, Entscheidung und gespeichertem
Ergebnis steht in Abschnitt~\ref{sec:eval-fortschreibung}.
 Diese fachliche Auswahl ist von der
Evidenzklärung in Abschnitt~\ref{sec:evidenzgebundene-verarbeitung}
zu unterscheiden; den Durchgang bis zu Core und Issues zeigt
Tabelle~\ref{t:beispiel-besuchsplanung}.

\emph{Andere Entscheidungsgegenstände.} Bei einem Widerspruch reicht
die Entscheidung über seine Aufnahme als offene Entscheidung (DEC)
noch nicht aus, um den Konflikt aufzulösen. Dafür stellt ein deterministischer
Schritt der Hauptkette offene Entscheidungen aus dem Core zusammen.
Der Mensch sieht den bisherigen Inhalt und die entgegenstehende
Aussage. Er kann das Original behalten, die neue Aussage übernehmen,
eine verfeinerte Fassung festlegen oder die Entscheidung vertagen.
Bei Übernahme oder Verfeinerung gibt er den endgültigen Text an;
der Workflow überträgt seine Auswahl in einen ausführbaren Plan.
Die möglichen Bestandsänderungen erklärt
Abschnitt~\ref{sec:core-projektfortschreibung}; Fall F3 in
Abschnitt~\ref{sec:eval-fortschreibung} zeigt eine tatsächlich
ausgeführte Ablösung. Am zentralen DEC-Gate bestimmt kein weiterer
Agent die Auflösung. Eine zusätzliche Begründung ist bei solchen
zielbezogenen Konflikten optional. Offene Fragen ohne Zielanforderung
haben eingeschränkte Optionen: Ihre Erledigung und der bewusste
Verzicht auf eine Architekturfestlegung benötigen eine Begründung.

Weitere Freigaben betreffen die Arbeitsplanung und die Außenwirkung.
Bei der PBI-Fortschreibung entscheidet der Mensch getrennt über
vorgeschlagene Zuordnungen und die inhaltliche Angleichung der
Arbeitspakete. Vor der GitHub-Projektion werden die extern wirksamen
Operationen gesondert freigegeben. Die vorgelagerte
Evidenz-Adjudikation klärt dagegen die verwendbaren Aussagen und
ihre Einordnung; sie autorisiert keine fachliche Core-Änderung.
Anhang~\ref{anh:gate-matrix} ordnet die zehn Entscheidungspunkte
der Hauptkette nach Gegenstand und Wirkung ein.

Im interaktiven Fortschreibungsbetrieb autorisiert weder ein Agent
noch eine deterministische Komponente den eigenen fachlichen
Vorschlag. Deklarierte Experimentbetriebsarten können die
Interaktion für definierte Untersuchungen durch vorgegebene oder
gespeicherte Entscheidungen ersetzen. Daraus entsteht keine
zusätzliche fachliche Autorität eines Agenten. Die technische
Realisierung und ihre Betriebsarten erklärt
Abschnitt~\ref{sec:realisierung-human-gates};
die zugehörigen Bedienwege beschreibt
Abschnitt~\ref{sec:realisierung-steward}.

Eine menschliche Freigabe garantiert keine fachliche Richtigkeit;
die Invariantenprüfung sichert keine semantisch richtige Beziehung.
Die Kontrollen begrenzen, wie Vorschläge wirksam werden.
Kapitel~\ref{chap:evaluation} untersucht ausgewählte Erzeugungs-,
Prüf- und Kontrollmechanismen. Wie der freigegebene Projektzustand
nach außen dargestellt und externe Information zurückgeführt wird,
erklärt der folgende Abschnitt.
